\section{Discussion and Limitations}
\label{sec:discussion}

\subsection{What CCS is — and is not}
\label{sec:discussion:what_ccs_is}

CCS is a \emph{measurement}, not a remedy. It tells the deployer that
a retrieved set is internally fragmented relative to what the
generator needs. It does not tell the deployer how to fix it. In our
deployment policy, the remediation is to \emph{not} refine and trust
the original retrieval. In other deployments the right move might be
to re-retrieve with a larger $k$, fall back to zero-shot, or trigger
external search. The diagnostic is decoupled from the response, by
design.

\subsection{Three regimes where CCS does not predict faithfulness}
\label{sec:discussion:where_ccs_fails}

A theory is more credible when its proponents articulate where it
breaks. We find three regimes in which CCS does not reliably predict
the paradox:

\paragraph{(1) Highly redundant corpora.}
When the corpus is dominated by near-duplicate passages (legal
boilerplate, software changelogs), every retrieval set has high CCS
by construction. The signal saturates and provides no
discrimination. Diagnostic: if the median CCS across queries exceeds
$0.85$, treat CCS as an availability check rather than a quality
signal.

\paragraph{(2) Long-context generators with single-document inputs.}
If the generator's context window comfortably holds the entire
relevant document and retrieval reduces to ``select this one
document'', coherence is trivially $1.0$ and the paradox does not
arise --- there is no fragmentation to begin with. This is consistent
with our long-form non-result on QASPER (Table~\ref{tab:robustness_summary}).

\paragraph{(3) Domain-mismatched encoders.}
When the encoder is general-purpose and the corpus is specialised
(MiniLM on biomedical PubMedQA), every retrieval set has \emph{low}
CCS because the encoder cannot resolve fine-grained in-domain
similarity. The CCS distribution collapses and the paradox is
attenuated --- but the failure mode is different: the issue is now
upstream encoder mismatch rather than refinement-induced
fragmentation. CCS still detects the regime (its distribution flags
the encoder mismatch) but cannot fix it.

The common thread: CCS measures a specific structural property of the
retrieved set. When that property is uninformative (saturated, trivial,
or collapsed), CCS is uninformative. When it is informative, it is a
clean predictor of refinement-induced failure.

\subsection{Limitations}
\label{sec:discussion:limits}

We mark four limitations explicitly:

\paragraph{Long-form generation.} The paradox is a property of
short-answer QA where the generator must extract a precise span. On
long-form generation (QASPER, MS-MARCO multi-sentence) the same
intervention reduces unsupported-claim rate but the paradox itself
does not generalise in the form measured by single-span NLI.

\paragraph{Magnitude vs.\ generic noise.} On SQuAD the noise sensitivity
slope ($-0.154$) exceeds the paradox magnitude ($0.100$). We can
claim the paradox is sign-distinct from generic retrieval noise on
multi-hop and biomedical datasets, but we cannot claim its magnitude
exceeds noise on SQuAD specifically. The original $\textrm{ratio} \geq 2$
target was not met; we report the result honestly.

\paragraph{Adversarial set yield.} The auto-generated $200$-case
adversarial set was filtered to $129$ validated cases by an NLI
validator + length/overlap heuristic. The mechanistic classifier in
the appendix is trained on a $20$-pair subset; full mechanistic
validation across the $129$-case set remains a follow-up.

\paragraph{Closed-model frontier.} We evaluate up to GPT-OSS-$120$B
via Groq's free tier; we do not evaluate $400$B+ closed models
(GPT-$4$, Claude-$3.5$). The coherence account predicts the paradox
should attenuate further at very large scale on tasks where
parametric memory can substitute for retrieval, and persist on tasks
where retrieval supplies novel information. This is a clean
prediction for follow-up work.
